A trustworthy explanation of a dashcam collision must name the right vehicles, state the right cause, and place the impact at the right moment.
We study whether Video-LLMs can do this on the Car Crash Dataset (CCD), and we separate two distinct ways they fail: \emph{omission} (leaving out what happened) and \emph{hallucination} (asserting what did not).
Four open Video-LLMs evaluated zero-shot on a 150-video held-out split are fluent but nearly useless: explanation BLEU-4 $\leq 0.026$, and $28$--$98\%$ of forensic fields are missing.
QLoRA adaptation of Qwen2.5-VL-7B on 1{,}198 training videos (\crashlogic) raises BLEU-4 to $0.142$ and BERTScore from $0.486$ to $0.686$, and cuts the omission cost from $0.462$ to $0.107$ ($p<10^{-4}$, improved on 140 of 150 videos).
The hallucination cost, however, does not move ($0.227\to0.223$, $p=0.81$): one in five adapted explanations names a vehicle whose colour contradicts the ground truth, and the model fabricates a collision in every confirmed no-crash test video.
We further show that the adapted model's apparently strong crash-window localisation (tIoU $0.012\to0.373$) is a learned dataset prior: it emits only two windows, its predicted start is uncorrelated with the true start (Spearman $\rho{=}0.009$), and a constant window beats it (tIoU $0.408$).
Finally, we test whether decode-time Temporal Contrastive Decoding (\tcd), which subtracts logits obtained from a reversed-frame negative, repairs the residual errors.
Under paired tests on 150 videos it does not: mild contrast ($\alpha{=}0.5$) is harmless but changes no metric significantly, while strong contrast ($\alpha{\geq}1.5$) significantly increases both hallucination and omission.
We explain why with a simple argument---reversal preserves every order-invariant cue (colour, weather, agent identity), so the contrast cannot touch the errors that dominate---and we release outputs, per-video scores, and an automatic error taxonomy so that future claims about faithful accident explanation can be checked against the same bar.
